% chapters/07/section-04.tex
\section{Kapan Persegi Pembatas Tidak Cukup}
\label{sec:mask}

Persegi pembatas dipilih karena murah dihitung dan cukup untuk sebagian besar
keperluan. Bagian ini membahas keterbatasannya, cara mengatasinya, dan kapan
cara itu memang diperlukan.

\subsection{Keterbatasan Persegi Pembatas}

Persegi pembatas selalu lebih besar daripada gambarnya, kecuali untuk benda
yang bentuknya memang persegi. Pada benda yang miring atau berlubang,
selisihnya dapat cukup besar sehingga dua benda dinilai bersentuhan padahal
gambarnya masih berjarak. Keadaan itu diperlihatkan pada
Gambar~\ref{fig:aabb-mask}.

\begin{figure}[htbp]
\centering
% chapters/07/figures/aabb-mask.tikz
% Kapan persegi pembatas tidak cukup: dua benda miring yang persegi
% pembatasnya beririsan, tetapi gambarnya tidak bersinggungan.
\begin{tikzpicture}[
  kotakpembatas/.style={draw=ubRed, thick, dashed},
  bentuk/.style={draw=ubPrimary, thick, fill=ubPrimary!20},
  panah/.style={-{Stealth[length=2.2mm]}, thick, draw=ubAccent},
  judul/.style={font=\footnotesize\bfseries, ubPrimary},
  label/.style={font=\scriptsize, align=center}
]

% ---------- Persegi pembatas beririsan ----------
\begin{scope}
  \node[judul, anchor=south] at (1.9,2.9) {diperiksa dengan persegi pembatas};

  \draw[bentuk] (0.35,0.35) -- (2.35,1.55) -- (0.35,1.55) -- cycle;
  \draw[bentuk] (1.45,0.35) -- (3.45,1.55) -- (1.45,1.55) -- cycle;
  \draw[kotakpembatas] (0.35,0.35) rectangle (2.35,1.55);
  \draw[kotakpembatas] (1.45,0.35) rectangle (3.45,1.55);

  \node[label, anchor=north, text width=52mm] at (1.9,-0.15)
    {Kedua persegi beririsan, sehingga program menyimpulkan kedua benda
     bersentuhan, padahal gambarnya tidak bersinggungan.};
\end{scope}

% ---------- Diperiksa dengan mask ----------
\begin{scope}[xshift=7.6cm]
  \node[judul, anchor=south] at (1.9,2.9) {diperiksa dengan \istilah{mask}};

  \draw[bentuk] (0.35,0.35) -- (2.35,1.55) -- (0.35,1.55) -- cycle;
  \draw[bentuk] (1.45,0.35) -- (3.45,1.55) -- (1.45,1.55) -- cycle;

  \node[label, anchor=north, text width=52mm] at (1.9,-0.15)
    {Yang dibandingkan adalah titik-titik yang benar-benar terisi. Karena tidak
     ada yang bertumpang tindih, kedua benda dinilai tidak bersentuhan.};
\end{scope}
\end{tikzpicture}

\caption{Persegi pembatas beririsan walaupun gambarnya tidak bersinggungan}
\label{fig:aabb-mask}
\end{figure}

Gejala yang muncul pada permainan mudah dikenali: pemain merasa ``kena'' oleh
sesuatu yang sebenarnya belum menyentuhnya. Pada permainan yang mengandalkan
ketepatan, misalnya permainan yang mengharuskan pemain melewati celah sempit,
selisih beberapa piksel itu menentukan berhasil atau gagalnya percobaan.

\subsection{Pemeriksaan dengan \istilah{Mask}}

Pygame menyediakan cara yang lebih cermat melalui \code{pygame.mask}. Sebuah
\istilah{mask} adalah gambaran hitam-putih yang menyatakan titik mana saja
pada sebuah permukaan yang benar-benar terisi.

\begin{kodekeluaran}
mask_pemain = pygame.mask.from_surface(player.image)
mask_bongkah = pygame.mask.from_surface(platform.image)

if mask_pemain.overlap(mask_bongkah, (dx, dy)):
    # ada titik yang benar-benar bertumpang tindih
\end{kodekeluaran}

Perhatikan argumen kedua pada \code{overlap}: nilainya adalah selisih letak
kedua benda, bukan letak salah satunya. Nilai itu diperoleh dari selisih
kedua \code{rect}, dan kesalahan menentukan nilai tersebut merupakan sumber
kekeliruan yang paling sering terjadi pada pemakaian \istilah{mask}.

Mask juga dapat dibuat sekali saja, yaitu ketika benda dibuat, lalu disimpan
pada benda itu. Membuat mask berulang kali setiap bingkai membuang waktu,
sebab pembuatannya menuntut pemeriksaan seluruh titik gambar.

\subsection{Memilih Metode yang Tepat}

Pilihan antara persegi pembatas dan \istilah{mask} adalah pertukaran antara
kecepatan perhitungan dan ketelitian. Tabel~\ref{tab:pilihan-metode}
merangkum pertimbangannya.

\begin{table}[htbp]
\centering
\caption{Pertimbangan memilih metode pemeriksaan tabrakan}
\label{tab:pilihan-metode}
{\setlength{\tabcolsep}{4pt}\footnotesize%
\begin{tabularx}{\textwidth}{@{}L{3.6cm} L{4.6cm} Y@{}}
\toprule
\headrow \textbf{Keadaan} & \textbf{Metode yang dianjurkan} & \textbf{Alasan} \\
\midrule
Benda berbentuk persegi atau hampir persegi &
Persegi pembatas &
Selisih dengan bentuk sebenarnya kecil, sedangkan perhitungannya paling
murah. \\
\addlinespace
Benda banyak dan bergerak cepat &
Persegi pembatas, dilengkapi penyelesaian per sumbu &
Ketelitian beberapa piksel tidak terasa pada benda yang bergerak cepat,
sedangkan \istilah{mask} memperberat setiap bingkai. \\
\addlinespace
Benda miring, berlubang, atau bergerigi &
\istilah{Mask} &
Persegi pembatasnya jauh lebih besar daripada gambar, sehingga tabrakan
palsu sering terjadi. \\
\addlinespace
Permainan yang menuntut ketepatan tinggi &
\istilah{Mask}, setelah pemeriksaan persegi pembatas &
Persegi pembatas dipakai lebih dahulu untuk menyaring, lalu \istilah{mask}
memeriksa hanya pada benda yang hampir bersentuhan. \\
\bottomrule
\end{tabularx}}
\end{table}

\begin{tipsbox}
Baris terakhir pada tabel tersebut memuat cara yang paling sering dipakai pada
permainan sungguhan: kedua metode digabungkan. Persegi pembatas diperiksa
lebih dahulu karena murah; hanya apabila hasilnya beririsan, \istilah{mask}
diperiksa. Dengan cara itu, \istilah{mask} yang mahal hanya dijalankan pada
beberapa benda yang memang berdekatan, bukan pada seluruh benda setiap
bingkai.
\end{tipsbox}

\begin{catatabox}
Pada program bab ini, seluruh benda berbentuk persegi sehingga pemeriksaan
persegi pembatas sudah memadai. Pilihan itu diambil secara sadar, sesuai
indikator \ind{a} pada bab ini: memilih metode yang \textbf{tepat}, bukan yang
paling canggih. Menambahkan \istilah{mask} pada program ini hanya akan
memperberat perhitungan tanpa mengubah hasilnya.
\end{catatabox}
